\SourcePage{149}{154}
\section{Motion in a centrally symmetric field}

Consider the motion of an electron in a centrally symmetric electric field. Since angular momentum and parity are conserved in motion in a central field (relative to the centre of the field, chosen as the coordinate origin), everything said in~\S\,24 about the angular dependence of wave functions of spherical waves of free particles applies to this motion as well. Only the radial functions change. Accordingly we seek the wave function of stationary states (in the standard representation) in the form
\begin{equation}
\varphi = \begin{pmatrix} \varphi \\ \chi \end{pmatrix}
= \begin{pmatrix} f(r)\,\Omega_{jlm} \\[4pt] (-1)^{(1+l-l')/2}\,g(r)\,\Omega_{jl'm} \end{pmatrix},
\tag{35.1}
\end{equation}
where $l = j \pm \tfrac{1}{2}$, $l' = 2j - l$, and the factor $-1$ is introduced to simplify the subsequent formulae.

The Dirac equation in the standard representation gives the following system of equations for $\varphi$ and $\chi$:
\begin{equation}
(\varepsilon - m - U)\varphi = \boldsymbol{\sigma}\widehat{\mathbf{p}}\,\chi,\qquad
(\varepsilon + m - U)\chi = \boldsymbol{\sigma}\widehat{\mathbf{p}}\,\varphi,
\tag{35.2}
\end{equation}
where $U(r) = e\Phi(r)$ is the potential energy of the electron in the field. The calculation of the result of substituting the expressions~(35.1) into these equations reduces to computing the right-hand sides.

Expressing the spherical spinor $\Omega_{jl'm}$ through $\Omega_{jlm}$ according to
\[
\Omega_{jl'm} = i^{l-l'}\!\left(\frac{\boldsymbol{\sigma}\cdot\mathbf{r}}{r}\right)\!\Omega_{jlm}
\]
(see~(24.8)), we write:
\[
(\boldsymbol{\sigma}\widehat{\mathbf{p}})\chi = -i(\boldsymbol{\sigma}\nabla)(\boldsymbol{\sigma}\mathbf{r})\,\frac{g}{r}\,\Omega_{jlm}.
\]

\SourcePage{150}{155}
Transforming now the product $(\boldsymbol{\sigma}\widehat{\mathbf{p}})(\boldsymbol{\sigma}\mathbf{r})$ with the aid of formula~(33.5), after expanding the vector operations we find
\[
(\boldsymbol{\sigma}\widehat{\mathbf{p}})\chi = -i\{\widehat p_r + i\boldsymbol{\sigma}[\widehat{\mathbf{p}}\,\widehat{\mathbf{r}}]\}\,\frac{g}{r}\,\Omega_{jlm} =
\]
\[
= \bigl\{-\mathop{\mathrm{div}}\mathbf{r} - (\mathbf{r}\nabla) - \boldsymbol{\sigma}[\mathbf{r}\widehat{\mathbf{r}}]\bigr\}\,\frac{g}{r}\,\Omega_{jlm}
= -\!\left\{g' + \frac{2}{r}g + \frac{g}{r}\,\boldsymbol{\sigma}\cdot\widehat{\mathbf{l}}\right\}\!\Omega_{jlm},
\]
where $\widehat{\mathbf{l}} = [\mathbf{r}\widehat{\mathbf{p}}]$ is the orbital angular momentum operator; the prime denotes differentiation with respect to~$r$. The eigenvalues of the product $\boldsymbol{\sigma}\cdot\widehat{\mathbf{l}} = 2\widehat{\mathbf{l}}\cdot\widehat{\mathbf{s}}$ are
\[
2\mathbf{l}\cdot\mathbf{s} = \mathbf{j}^2 - \mathbf{l}^2 - \mathbf{s}^2 = j(j+1) - l(l+1) - \tfrac{3}{4} =
\begin{cases}
j - \tfrac{1}{2}, & l = j - \tfrac{1}{2},\\
-j - \tfrac{3}{2}, & l = j + \tfrac{1}{2}.
\end{cases}
\]

For uniformity of notation in both cases ($l = j \pm \tfrac{1}{2}$) it is convenient to introduce the notation
\begin{equation}
\varkappa =
\begin{cases}
-(j + \tfrac{1}{2}) = -(l+1), & j = l + \tfrac{1}{2},\\
+(j + \tfrac{1}{2}) = l, & j = l - \tfrac{1}{2}.
\end{cases}
\tag{35.3}
\end{equation}
The number $\varkappa$ takes all integer values excluding zero (positive numbers correspond to $j = l - \tfrac{1}{2}$, and negative ones to $j = l + \tfrac{1}{2}$). Then $\mathbf{l}\cdot\boldsymbol{\sigma} = -(1 + \varkappa)$, so that
\[
(\boldsymbol{\sigma}\widehat{\mathbf{p}})\chi = -\!\left(g' + \frac{1-\varkappa}{r}\,g\right)\!\Omega_{jlm}.
\]

On substituting this expression into the first of equations~(35.2), the spherical spinor $\Omega_{jlm}$ cancels on both sides of the equation. Proceeding analogously with the second equation, we obtain the following system for the radial functions:
\begin{equation}
\begin{aligned}
f' + \frac{1+\varkappa}{r}\,f - (\varepsilon + m - U)\,g &= 0,\\[4pt]
g' + \frac{1-\varkappa}{r}\,g + (\varepsilon - m - U)\,f &= 0,
\end{aligned}
\tag{35.4}
\end{equation}
or
\begin{equation}
\begin{aligned}
(fr)' + \frac{\varkappa}{r}\,(fr) - (\varepsilon + m - U)\,gr &= 0,\\[4pt]
(gr)' - \frac{\varkappa}{r}\,(gr) + (\varepsilon - m - U)\,fr &= 0.
\end{aligned}
\tag{35.5}
\end{equation}

\SourcePage{151}{156}
Let us investigate the behaviour of $f$ and $g$ at small distances, assuming that the field $U(r)$ grows as $r\to 0$ faster than $1/r$. Then in the region of small $r$ equations~(35.4) take the form
\[
f' + Ug = 0,\qquad g' - Uf = 0.
\]
They have real solutions of the form
\begin{equation}
f = \mathrm{const}\cdot\sin\!\left(\int\!U\,dr + \delta\right),\qquad
g = \mathrm{const}\cdot\cos\!\left(\int\!U\,dr + \delta\right),
\tag{35.6}
\end{equation}
where $\delta$ is an arbitrary constant. These functions oscillate as $r\to 0$, without tending to any limit. It is easy to see that this situation corresponds in the non-relativistic theory to ``falling'' of the particle to the centre.

First of all we note that the region of small distances imposes in this case no restrictions on the choice of solution: the condition at $r = 0$ for the oscillating function is absent, and the choice of the constant $\delta$ remains arbitrary (the correct behaviour of the wave function at large $r$ can be achieved for any $\varepsilon$ by an appropriate choice of $\delta$). This indeterminacy can be removed by considering the singular (at $r = 0$) potential as the limit $r_0\to 0$ of the potential ``cut~off'' at some small~$r_0$ (i.e.\ equal to $U(r)$ for $r > r_0$ and $U(r_0)$ for $r < r_0$). With a finite $r_0$ one obtains, of course, a definite system of energy levels. However, the energy of the ground state tends to $-\infty$ as $r_0\to 0$.

In the non-relativistic theory this is precisely the ``falling'' to the centre, since a particle at a deep level is localised in a small region around $r = 0$. In the relativistic theory, however, such a situation is altogether inadmissible, since it signifies the instability of the system with respect to the spontaneous creation of electron--positron pairs. Indeed, if in vacuum the energy needed to create such a pair exceeds~$2m$, in the field a smaller energy already suffices. If there exists a bound state of the electron with energy $\varepsilon < m$, pair creation is possible at the cost of only the energy $\varepsilon + m < 2m$, with a free positron created and the electron in the bound state. If the energy of the bound state $\varepsilon < -m$, such a field can create positrons (with energy $-\varepsilon > m$) spontaneously, without any input of energy from an external source. In the field being considered, as $r_0\to 0$ there is an infinite number of such ``anomalous'' levels with $\varepsilon < -m$. Therefore a potential $\Phi(r)$ growing as $r\to 0$ faster than $1/r$ cannot at all be considered in the Dirac theory. Let us stress that this applies to potentials of either sign. The ``falling'' occurs, of course, only in the case of attraction, since the sign of $U = e\Phi$ depends also on the sign of the charge; in one case the electronic, and in the other the positronic levels behave anomalously; in the second case the field creates free electrons.

\SourcePage{152}{157}
Let us consider further the behaviour of wave functions at large distances. If the field $U(r)$ falls off sufficiently rapidly at $r\to\infty$, then in determining the asymptotic form of the wave functions at large distances the field can be completely neglected in the equations. For $\varepsilon > m$, i.e.\ in the continuous-spectrum region, we return then to the free-motion equation, so that the asymptotic form of the wave functions (spherical waves) differs from that for a free particle only by the appearance of additional ``phase shifts'', whose values are determined by the form of the field at close distances\,\footnote{Cf.\ III, \S\,33. As in the non-relativistic theory, $U(r)$ must fall off faster than $1/r$. The case $U \sim 1/r$ will be considered separately in~\S\,36.}. These shifts depend on the values of $j$ and~$l$, or, equivalently, on the number $\varkappa$ (and also, of course, on the energy~$\varepsilon$). Denoting them by $\delta_\varkappa$ and using the expression for the free spherical wave~(24.7), we can immediately write the desired asymptotic formula
\begin{equation}
\psi \approx \frac{2}{r}\,\frac{1}{\sqrt{2\varepsilon}}
\begin{pmatrix}
\sqrt{\varepsilon+m}\;\Omega_{jlm}\sin\!\left(pr - \dfrac{\pi l}{2} + \delta_\varkappa\right)\\[6pt]
-\sqrt{\varepsilon-m}\;\Omega_{jl'm}\sin\!\left(pr - \dfrac{\pi l'}{2} + \delta_\varkappa\right)
\end{pmatrix},
\tag{35.7}
\end{equation}
or, taking into account the definition~(35.1):
\begin{equation}
\left.\begin{aligned}
f\\g
\end{aligned}\right\}
= \frac{\sqrt{2}}{r}\sqrt{\frac{\varepsilon\pm m}{\varepsilon}}\;\sin\!\left(pr - \frac{\pi l}{2} + \delta_\varkappa\right),
\tag{35.8}
\end{equation}
where $p = \sqrt{\varepsilon^2 - m^2}$. The overall coefficient here corresponds to normalisation of the radial functions according to~(24.5).

The wave functions of the discrete spectrum ($\varepsilon < m$) decay exponentially at $r\to\infty$ according to the law
\begin{equation}
f = -\sqrt{\frac{m+\varepsilon}{m-\varepsilon}}\;g = \frac{A_0}{r}\exp\!\left(-r\sqrt{m^2 - \varepsilon^2}\right),
\tag{35.9}
\end{equation}
where $A_0$ is a constant.

As in the non-relativistic theory, the phase shifts $\delta_\varkappa$ (more precisely, the quantities $e^{2i\delta_\varkappa} - 1$) determine the scattering amplitude in the given field (this will be discussed in more detail in~\S\,37). We shall not investigate here the analytic properties of these quantities (cf.\ III, \S\,128). We note only that $e^{2i\delta_\varkappa}$ as a function of the energy still has poles at the points corresponding to the bound-state energy levels of the particle. The residue of the function $e^{2i\delta_\varkappa}$ at such a pole is related in a definite way to the coefficient in the asymptotic expression of the corresponding wave function of the discrete spectrum. Let us find this relation, which generalises the non-relativistic formula~(128.17) (see~III). The necessary calculations are entirely analogous to those performed in~III, \S\,128.

\SourcePage{153}{158}
Differentiate equations~(35.5) with respect to the energy:
\[
\left(\frac{\partial\,rf}{\partial\varepsilon}\right)' + \frac{\varkappa}{r}\,\frac{\partial\,rf}{\partial\varepsilon} - (\varepsilon + m - U)\,\frac{\partial\,rg}{\partial\varepsilon} = rg,
\]
\[
\left(\frac{\partial\,rg}{\partial\varepsilon}\right)' - \frac{\varkappa}{r}\,\frac{\partial\,rg}{\partial\varepsilon} + (\varepsilon - m - U)\,\frac{\partial\,rf}{\partial\varepsilon} = -rf.
\]
Multiply these two equations respectively by $rg$ and $-rf$, and the two equations~(35.5) respectively by $-rg$ and $rf$; then add all four equations term by term. After all cancellations we get
\[
\frac{\partial}{\partial r}\!\left[r^2\!\left(g\,\frac{\partial f}{\partial\varepsilon} - f\,\frac{\partial g}{\partial\varepsilon}\right)\right] = r^2(f^2 + g^2).
\]
Integrate this relation over $r$:
\[
r^2\!\left(g\,\frac{\partial f}{\partial\varepsilon} - f\,\frac{\partial g}{\partial\varepsilon}\right) = \int_0^r (f^2 + g^2)\,r^2\,dr,
\]
and then pass to the limit $r\to\infty$. By the normalisation condition the integral on the right-hand side becomes unity. On the left-hand side we take into account that in the asymptotic region the functions $f$ and $g$ are related by
\[
rg = \frac{(rf)'}{\varepsilon + m},
\]
obtained from~(35.5) by neglecting the terms with $U$ and $1/r$. As a result we get
\begin{equation}
\frac{1}{\varepsilon + m}\!\left[(rf)'\,\frac{\partial\,rf}{\partial\varepsilon} - rf\!\left(\frac{\partial\,rf}{\partial\varepsilon}\right)'\right] = 1.
\tag{35.10}
\end{equation}

This formula differs from the analogous non-relativistic formula (for the function $\chi$) only by the coefficient ($\varepsilon + m$ instead of $2m$). There is therefore no need to repeat all the subsequent calculations, and we immediately give the final formula, valid near the point $\varepsilon = \varepsilon_0$ ($\varepsilon_0$ being the energy level):
\begin{equation}
e^{2i\delta_\varkappa} = (-1)^l\,\frac{2A_0^2}{\varepsilon - \varepsilon_0}\,\sqrt{\frac{m - \varepsilon_0}{m + \varepsilon_0}},
\tag{35.11}
\end{equation}
where $A_0$ is the coefficient in the asymptotic expression~(35.9).

\SourcePage{154}{159}
\begin{center}
\textbf{Problem}
\end{center}

Find the limiting form of the wave function at small $r$ in the field $U \sim r^{-s}$, $s < 1$.

\textbf{Solution.} For a free particle we have at small~$r$: $f \sim r^l$, $g \sim r^{l'}$, so that for $l < l'$: $f \gg g$, and for $l > l'$: $f \ll g$. We make the assumption (justified by the result) that such a relation is preserved also in the field being considered. For $l < l'$ (i.e.\ $l = j - \tfrac{1}{2}$, $\varkappa = -l-1$) the $g$ term in the first of equations~(35.4) can be dropped, so that $f \sim r^l$ as before. From the second equation we then have $g \sim rfU$, i.e.\ $g \sim r^{l+1-s}$. The case $l > l'$ is treated analogously. As a result we find:
\[
\begin{aligned}
\text{for } l < l'\!: \quad f &\sim r^l, & g &\sim r^{l'-s};\\
\text{for } l > l'\!: \quad f &\sim r^{l-s}, & g &\sim r^{l'}.
\end{aligned}
\]
